\documentclass[licencjacka, en]{pracamgr}

% Dane magistrantów:
\autor{Heorhii Lopatin}{456366}

\title{The Deep Ritz Method }
\titlepl{Głęboka metoda Ritz'a}

\usepackage{hyperref}
\usepackage{graphicx}
\usepackage{placeins}
\usepackage{subcaption}
\usepackage{amssymb}
\usepackage{amsmath}
\usepackage{xurl}

\kierunek{Mathematics}

\opiekun{dr Piotr Krzyżanowski\\
  Institute of Applied Mathematics and Mechanics\\
  }

\date{June 2024}

\dziedzina{ 
% 11.4 Artificial Intelligence\\ 
}

\klasyfikacja{}


% Słowa kluczowe:
\keywords{}

\newtheorem{defi}{Definicja}[section]
\newcommand{\R}{\mathbb{R}}

% Our own chapter definition
\begin{document}
\newcommand{\rozdzial}[2]{
	\chapter{#1}
	\label{chap:#2}
	\input{chapters/#2}
}

\maketitle

%tu idzie streszczenie na strone poczatkowa
\begin{abstract}
	% This thesis concerns research into the use of machine learning
	% and large language models in market analysis, focusing on market
	% predictions.
\end{abstract}

\tableofcontents

\rozdzial{Introduction}{introduction}


\begin{thebibliography}{99}
	\addcontentsline{toc}{chapter}{Bibliography}

\end{thebibliography}

\end{document}


%%% Local Variables:
%%% mode: latex
%%% TeX-master: t
%%% coding: latin-2
%%% End:
